\documentclass[10pt,a4paper]{amsart}
\usepackage[margin=19mm]{geometry}
\usepackage{amsmath,amssymb,mathtools}
\usepackage[scr=rsfs,cal=euler]{mathalfa}
\usepackage{xfrac}
\usepackage[unicode,hidelinks,bookmarks=false]{hyperref}
\newcommand{\ie}{i.e.\ }
\newcommand{\cf}{cf.\ }
\newcommand{\resp}{respectively\ }
\newcommand{\Subsection}{\subsection}
\providecommand{\nobreakdash}{\nobreak\hbox{-}}
\setlength{\emergencystretch}{2em}
\multlinegap0pt
\usepackage{bm}
\usepackage{bbm}
\usepackage{dsfont}
\usepackage{xspace}
\usepackage{cancel}
\usepackage{mathtools}
\usepackage{amsthm}
\makeatletter
\@ifundefined{theorem}{\newtheorem{theorem}{Theorem}}{}
\makeatother
\title[Krall-type orthogonal polynomials and integrable isomonodrom]{Krall-type orthogonal polynomials and integrable isomonodromic deformations}
\author{Luc Haine}
\date{Source: arXiv:2603.01132v2 (CC BY 4.0)}
\begin{document}
\maketitle

\noindent\emph{Source and licence.} Krall-type orthogonal polynomials and integrable isomonodromic deformations. Version arXiv:2603.01132v2 (CC BY 4.0), distributed under the Creative Commons Attribution 4.0 International licence, as recorded in the arXiv licence metadata for this exact version and shown on the abstract page https://arxiv.org/abs/2603.01132v2. Full rights notice: licenses/arxiv-2603.01132-CC-BY-4.0.txt.\par\smallskip

\noindent\textit{Editorial scope.} Counted unit: the Theorem whose source label is \texttt{theorem 5.2} in the article (source0.tex lines 1302--1322, source label \texttt{theorem 5.2}), delivered together with the article's own complete original proof of it (source0.tex lines 1323--1392, one \texttt{proof} environment of 70 lines). No mathematics is written, repaired or re-derived: statement and proof are the article's own text line for line. Required context retained uncounted: Kga (lines 963--965); PK1 (lines 982--985); aKy2 (lines 987--990); aKz2 (lines 987--990); ab (lines 992--994); zy (lines 1068--1070); rs (lines 1245--1247); SKFI1 (lines 1249--1252); SKFI2 (lines 1249--1252); drt1 (lines 1269--1271); drt2 (lines 1278--1280). Every definition, display and label of those lines is retained unchanged. Omitted from this package and not redistributed: 1--962, 966--981, 986--986, 991--991, 995--1067, 1071--1244, 1248--1248, 1253--1268, 1272--1277, 1281--1301, 1393--1893. Nothing inside a retained excerpt is omitted or altered: every retained body is delivered exactly as the article's lines are written, so no omission receipt of any kind accompanies this package. Citations: the retained excerpts carry no citation command at all, so no bibliography entry of the article is transcribed and no reference is redistributed. Reference adaptation: none was needed and none was made; every reference inside the delivered unit resolves inside it, so no unresolved reference or citation is produced. Printed slips kept literal: no printed word, symbol, hypothesis or number of the counted statement or of its proof was altered. Layout and class: the article's own class is \texttt{amsart} with packages including amsmath, amsfonts, amssymb, amsthm, eucal, url, hyperref, cleveref; the wrapper is the lane's pinned \texttt{amsart} editorial wrapper at 10pt on A4, which restates the theorem-like environments the retained text needs and adds the article's own macro definitions that the retained text needs (none), with extra wrapper packages (bm, bbm, dsfont, xspace, cancel, mathtools). The delivered numbering is the unit's own. Assets: no retained span contains a figure environment, a graphics-inclusion command, a PDF-page inclusion, a table or a TikZ picture; no image file is copied into this package, the package declares no asset dependency and no image file is redistributed with it. Licence: version arXiv:2603.01132v2 of this article is distributed under the Creative Commons Attribution 4.0 International licence, as recorded in the arXiv licence metadata for this exact version and shown on the abstract page https://arxiv.org/abs/2603.01132v2. A full rights notice is delivered at licenses/arxiv-2603.01132-CC-BY-4.0.txt. Characters: every retained excerpt is pure ASCII, so the delivered engine, XeLaTeX with its default Latin Modern text font, reports no missing character.
\begin{equation}
\gamma_n=\alpha+\beta+2n+1,\;\mbox{for}\; n=0,1,2,\ldots, \label{Kga}
\end{equation}

\begin{gather}
(\beta+1)t_1\frac{\partial y_n}{\partial t_1}=(\alpha+1)t_2\frac{\partial z_n}{\partial t_2},\label{PK1}\\
(\alpha+1)z_nt_2\frac{\partial y_n}{\partial t_2}+(\beta+1)y_nt_1\frac{\partial z_n}{\partial t_1}=(\beta+1)(y_n+z_n-4)t_1\frac{\partial y_n}{\partial t_1}-y_nz_n.\label{PK2}
\end{gather}

\begin{align}
y_n(t_1,t_2,\alpha,\beta)&=\frac{a}{t_2}+\frac{a^2(\beta^2-(\alpha+1)^2-\gamma_n^2)}{8(\alpha+1)^2t_2^2}-\frac{ab(\alpha+\beta+1)}{4 (\beta+1)t_1t_2}+O\Big(\Big(\frac{1}{t_1^2} +\frac{1}{t_2^2}\Big)^\frac{3}{2}\Big),\label{aKy2}\\
z_n(t_1,t_2,\alpha,\beta)&=\frac{b}{t_1}+\frac{b^2(\alpha^2-(\beta+1)^2-\gamma_n^2)}{8(\beta+1)^2t_1^2}-\frac{ab(\alpha+\beta+1)}{4(\alpha+1)t_1t_2}+O\Big(\Big(\frac{1}{t_1^2} +\frac{1}{t_2^2}\Big)^\frac{3}{2}\Big),\label{aKz2}
\end{align}

\begin{equation}
a=\frac{4(\alpha+1)(\alpha+1)_n(\alpha+\beta+2)_{n-1}}{n!(\beta+1)_n}, \quad b=\frac{4(\beta+1)(\beta+1)_n(\alpha+\beta+2)_{n-1}}{n!(\alpha+1)_n}.\label{ab}
\end{equation}

\begin{equation}
z_n(t_1,t_2,\alpha,\beta)=y_n(t_2,t_1,\beta,\alpha).\label{zy}
\end{equation}

\begin{equation} 
r_n=t_1\frac{\partial}{\partial t_1} \log\frac{y_n}{z_n},\quad s_n=\frac{t_1}{z_n}\frac{\partial}{\partial t_1}\log y_n,\label{rs}
\end{equation}

\begin{gather}
r_n^2-4z_ns_n^2+\frac{\gamma_n^2z_n}{4(\beta+1)^2}=1,\label{SKFI1}\\
\big((\beta+1)(4s_n-r_n)+1\big)^2-4(\beta+1)^2y_ns_n^2+\frac{\gamma_n^2 y_n}{4}=(\alpha+1)^2,\label{SKFI2}
\end{gather} 

\begin{equation}
t_1\frac{\partial r_n}{\partial t_1}=-2z_ns_n^2+\frac{\gamma_n^2 z_n}{8(\beta+1)^2}.\label{drt1}
\end{equation}

\begin{equation}
\frac{\partial r_n}{\partial t_2}=\frac{1}{2t_2(\alpha+1)(\beta+1)y_nz_n}\Big\{t_2^2(\alpha+1)^2\Big(\frac{\partial z_n}{\partial t_2}\Big)^2-t_1^2(\beta+1)^2\Big(\frac{\partial y_n}{\partial t_1}\Big)^2\Big\}=0.\label{drt2}
\end{equation}

\begin{theorem} \label{theorem 5.2}
The rational functions $y_n(t_1,t_2,\alpha,\beta)$ and $z_n(t_1,t_2,\alpha,\beta)$ that characterize the Koornwinder polynomials are explicitly given by
\begin{equation} 
y_n(t_1,t_2,\alpha,\beta)=\frac{4(\alpha+1)^2k_2t_2(k_1t_1+n)^2}{\sigma_n(t_1,t_2,\alpha,\beta)\tau_n(t_1,t_2,\alpha,\beta)},\; n\geq 0,\label{Ky}
\end{equation}
where 
\begin{align}
\sigma_n(t_1,t_2,\alpha,\beta)&=n^2(n-1)+n(\beta+n)k_1t_1+n(\alpha+n)k_2t_2
+(\alpha+\beta+n+1)k_1k_2t_1t_2,\nonumber\\
\tau_n(t_1,t_2,\alpha,\beta)&=k_1k_2t_1t_2+(\alpha+n+1)k_1t_1+(\beta+n+1)k_2t_2
+n(\alpha+\beta+n+2),\label{Ksigmatau}
\end{align}
and the coefficients $k_1, k_2$ are defined as follows:
\begin{equation} \label{k1k2}
k_1=\frac{n!(\beta+1)(\alpha+1)_n}{(\beta+1)_n(\alpha+\beta+2)_n},\; k_2=\frac{n!(\alpha+1)(\beta+1)_n}{(\alpha+1)_n(\alpha+\beta+2)_n}.
\end{equation}
Moreover, we have
\begin{equation}
z_n(t_1,t_2,\alpha,\beta)=y_n(t_2,t_1,\beta,\alpha).\label{Kz}
\end{equation}
\end{theorem}

\begin{proof}
The strategy of the proof is to first determine $r_n$ and $s_n$, from which the final result will follow. From \eqref{drt1}, using \eqref{SKFI1}, we obtain 
\begin{equation} 
2t_1\frac{\partial r_n}{\partial t_1}=1-r_n^2.\label{rt1}
\end{equation}
From \eqref{drt2}, $r_n$ depends only on $t_1$. Hence, integrating \eqref{rt1} gives
\begin{equation}
r_n=\frac{k_1t_1-n}{k_1t_1+n},\label{solr}
\end{equation}
with $k_1\neq 0$ being an arbitrary constant that depends on $n,\alpha, \beta$, still to be determined. We have chosen to write $r_n$ in this form so that the limit exists when $n=0$. Remembering the definition of $r_n$ \eqref{rs}, from the expansions \eqref{aKy2} and \eqref{aKz2}, we find that
\begin{equation} \label{k1}
k_1=\frac{4(\beta+1)^2}{(\alpha+\beta+n+1)b},
\end{equation} 
which, using \eqref{ab}, gives the expression for $k_1$ in \eqref{k1k2}. Defining
\begin{equation*}
\tilde{r}_n(t_1,t_2,\alpha,\beta)=r_n(t_2,t_1,\beta,\alpha),
\end{equation*}
it follows that $\tilde{r}_n$ depends only on $t_2$ and is given by
\begin{equation}\label{trt2}
\tilde{r}_n=\frac{k_2t_2-n}{k_2t_2+n},
\end{equation}
with $k_2$ obtained by permuting $\alpha$ with $\beta$ in $k_1$,
\begin{equation}\label{k2}
k_2=\frac{4(\alpha+1)^2}{(\alpha+\beta+n+1)a},
\end{equation}
thus yielding the expression for $k_2$ in \eqref{k1k2}, using \eqref{ab}.

Recalling again the definition of $r_n$ \eqref{rs}, we can integrate \eqref{solr} to obtain the quotient
\begin{equation}
q_n(t_1,t_2,\alpha,\beta)=\frac{y_n(t_1,t_2,\alpha,\beta)}{z_n(t_1,t_2,\alpha,\beta)}=f(t_2,\alpha,\beta)\frac{(k_1t_1+n)^2}{t_1},\label{q}
\end{equation}
where $f(t_2,\alpha,\beta)$ is an arbitrary function of $t_2$ depending on $\alpha,\beta$. Recalling \eqref{zy}, we have
\begin{equation*}
q_n(t_2,t_1,\beta,\alpha)=\frac{1}{q_n(t_1,t_2,\alpha,\beta)},
\end{equation*}
hence, we deduce that necessarily
\begin{equation*}
f(t_2,\alpha,\beta)=c(\alpha,\beta) \frac{t_2}{(k_2t_2+n)^2},
\end{equation*}
with $c(\alpha,\beta)$ a constant still to be determined. Setting $t_1=t_2=t$ in \eqref{q}, using \eqref{aKy2}, \eqref{aKz2} and taking the limit as $t\to\infty$, we find that
\begin{equation*}
c(\alpha,\beta)=\frac{ak_2^2}{bk_1^2}=\frac{(\alpha+1)^2k_2}{(\beta+1)^2k_1},
\end{equation*}
where the last equality is obtained from \eqref{k1} and \eqref{k2}. Putting everything together leads to
\begin{equation} 
q_n=\frac{(\alpha+1)^2k_2t_2(k_1t_1+n)^2}{(\beta+1)^2k_1t_1(k_2t_2+n)^2}=\frac{(\alpha+1)^2(\tilde{r}_n^2-1)}{(\beta+1)^2(r_n^2-1)}.\label{solq}
\end{equation} 

Solving \eqref{SKFI1} and \eqref{SKFI2} for $y_n$ and $z_n$ and taking the quotient, gives 
\begin{equation} \label{qSKFI}
\big((\beta+1)(4s_n-r_n)+1\big)^2-(\alpha+1)^2=(\beta+1)^2(r_n^2-1)q_n,
\end{equation}
which is a quadratic equation for $s_n$ in terms of $r_n$ and $q_n$. Substituting \eqref{solq} into \eqref{qSKFI}, a straightforward computation yields
\begin{equation} 
s_n=\pm \frac{(\alpha+1)\tilde{r}_n}{4(\beta+1)} +\frac{r_n}{4}-\frac{1}{4(\beta+1)}.\label{sols}
\end{equation}
From the definition of $s_n$ \eqref{rs}, using \eqref{PK1},we observe that
\begin{equation}
s_n(t_2,t_1,\beta,\alpha)=\frac{\beta+1}{\alpha+1} s_n(t_1,t_2,\alpha,\beta).\label{pstalbe} 
\end{equation}
We see at once that only the solution with the plus sign in \eqref{sols} satisfies \eqref{pstalbe}; hence, the other one is eliminated. Using \eqref{solq} and \eqref{qSKFI}, \eqref{SKFI2} can be written as
\begin{equation*}
(\alpha+1)^2(\tilde{r}_n^2-1)=\Big(4(\beta+1)^2s_n^2-\frac{\gamma_n^2}{4}\Big)y_n,
\end{equation*}
from which using \eqref{sols} with the plus sign, we obtain 
\begin{equation}
y_n=\frac{4(\alpha+1)^2(\tilde{r}_n^2-1)}{[(\alpha+1)\tilde{r}_n+(\beta+1)r_n+\gamma_n-1][(\alpha+1)\tilde{r}_n+(\beta+1)r_n-\gamma_n-1]}.\label{Kyrtr}
\end{equation}
Using \eqref{Kga}, \eqref{solr} and \eqref{trt2}, \eqref{Kyrtr} is equivalent to \eqref{Ky}. As we already know from \eqref{zy}, $z_n$ is obtained by permuting $t_1$ with $t_2$ and $\alpha$ with $\beta$ in \eqref{Kyrtr}, establishing \eqref{Kz}. Notice that $\sigma_n$ and $\tau_n$, as defined in \eqref{Ksigmatau}, are invariant under this permutation. This completes the proof of Theorem~\ref{theorem 5.2}.
\end{proof}

\end{document}
